\chapter*{요약과 문제의식}
\addcontentsline{toc}{chapter}{요약과 문제의식}
\markboth{요약과 문제의식}{}

\section*{한 쪽 요약}
\begin{keybox}[결론]
\begin{enumerate}[leftmargin=*]
\item \textbf{돈은 같은 등급 안에 있다.} LC 등급이 같아도 대출마다 질이 다르고, 그 차이를 고르면 아무거나 고를 때보다 매년 약 1\%p 더 번다(95\% 구간 0.95--1.11). 등급마다 상위 30\%를 사고 하위 30\%를 판다고 가정한 스프레드는 48개월 \textbf{전부} 양수였다(+2.38\%p).
\item \textbf{수익과 안정성은 다른 축이다.} 기대수익으로 고르면 어떤 기준(허들)에서도 가장 많이 번다. 안정성(과제 지표인 포트폴리오 Sharpe)의 1위는 투자자의 요구수익률에 따라 바뀐다. 국채 수준이면 부도 위험이 낮은 대출 위주가, 회사채(BAA) 수준이면 기대수익 순이 1위다. 국채 기준에서 보수적 선별이 이기는 것은 36개월 고위험 등급(E--G)이 위험만큼 금리를 받지 못했기 때문이다.
\item \textbf{투명한 규칙 하나가 둘 다를 잡을 가능성이 있다.} 등급마다 ``최근 24개월에 새 계좌를 적게 연'' 대출을 고르면, 네 가지 허들 모두에서 Sharpe가 가장 높고 수익은 기대수익 순과 거의 같았다. 다만 같은 데이터에서 보고 고른 규칙이라 test 전에 미리 고정한 \textbf{후보}다.
\item \textbf{재는 방법이 결론을 바꾼다.} 수익률 정의, 학습 시점, 채점 표본(끝난 대출만 보는지)을 바꾸면 같은 데이터의 결론이 뒤집혔다. 이 프로젝트의 가장 큰 기여는 모델이 아니라 \textbf{정직한 성적표}다.
\end{enumerate}
\Intuit{같은 가격표가 붙은 사과 중 더 좋은 사과를 고르면 더 번다. 비싼 사과(고위험 등급)는 가격만큼 좋지 않았다.}
\end{keybox}

\section*{문제의식: 왜 이 프로젝트가 어려운가}
과제는 ``Sharpe ratio를 최대로 만드는 대출 승인 전략''이다. 겉보기에는 부도 예측 문제다. 부도를 잘 맞히는 모델을 만들고, 부도할 것 같은 대출을 거절하면 될 것 같다. 우리는 이 직관에 숨어 있는 함정을 하나씩 확인했고, 그 과정이 곧 이 보고서의 뼈대다.

\begin{table}[H]\centering\small
\begin{tabularx}{\linewidth}{@{}p{38mm}LL@{}}\toprule
함정 & 왜 함정인가 & 우리가 한 일\\\midrule
부도만 맞히면 된다 & LC 금리는 이미 위험의 가격이다. 안전한 대출만 고르면 덜 번다 & 부도확률 $p$를 금리와 합친 \textbf{기대수익 점수}로 고른다(5장)\\
정확도(accuracy)로 평가한다 & 부도 비율이 약 15\%라, 전부 ``정상''이라 해도 85\%가 맞는다. 정확도는 돈과 무관하다 & 실제 현금흐름 수익률(IRR)로 채점한다(2장)\\
수익률은 아무렇게나 잰다 & 같은 부도 대출이 계산법에 따라 0\%, $-27\%$, $-77\%$가 되고, 결론이 뒤집힌다 & 수익률 정의 네 가지를 비교하고 IRR을 쓴다(2장)\\
데이터를 무작위로 나눠 시험한다 & 미래 대출로 과거를 맞혀 성적이 부풀려진다(AUC 0.73 vs 0.62--0.68) & 그때 알던 정보로만 학습한다(시점 학습, 3장)\\
끝난 대출만 보고 채점한다 & 빨리 끝난 대출(빨리 부도 난 대출)이 과대 대표되어 ``부도만 피하는'' 전략이 좋아 보인다 & 일부러 일찍 잘라 보는 시뮬레이션으로 편향을 측정하고 test 해석 규칙을 정했다(3장, 9장)\\
Sharpe 하나로 판단한다 & 대출 건별 Sharpe는 승인만 줄여도 변하고, 포트폴리오 Sharpe는 규모를 무시한다 & Sharpe 정의 세 가지를 모두 계산하고, 수익과 안정성 두 축으로 본다(6장)\\
비교 기준을 하나로 고정한다 & 국채 기준에서 맞는 결론이 회사채 기준에서 뒤집혔다 & 허들을 국채, +1\%p, +2\%p, BAA로 바꿔 가며 결론의 조건을 밝혔다(7장)\\
승인 대상을 자르는 단위를 신경 쓰지 않는다 & 연도 단위로 자르면 달마다 승인 금액이 들쭉날쭉해 Sharpe가 왜곡된다(외부 검증으로 발견) & 발행월 단위로 다시 계산했다(7, 8장)\\
좋은 결과는 믿는다 & 같은 데이터를 수십 번 보면 우연에 맞춰진다 & 판정 기준을 실행 전에 적고, 틀린 판단을 기록하고, test 규칙을 미리 고정했다(8, 9장)\\\bottomrule
\end{tabularx}
\caption{이 프로젝트가 넘어야 했던 함정들.}
\end{table}

\begin{warnbox}[이 보고서의 태도(11월 16일 개정)]
우리 결론 중 일부는 처음 판단과 반대로 바뀌었다. 예를 들어 ``Sharpe는 승인비율에 끌려가므로 주 지표로 부적합하다''는 팀원의 지적으로 범위를 좁혔고, ``LC 등급끼리의 가격은 공정했다''는 36개월 고위험 등급에서 틀렸다. 이런 정정은 숨기지 않고 8장에 모았다. 정정의 기록이 결과의 신뢰도를 만든다고 보기 때문이다.
\end{warnbox}

\section*{읽는 길잡이}
\begin{itemize}[leftmargin=*]
\item \textbf{1--4장(기초):} 과제와 데이터, 수익률을 재는 법, 시간을 지키는 법, EDA와 결측. 처음 보는 사람은 여기부터.
\item \textbf{5장:} 부도확률 $p$가 어떻게 만들어지고 \textbf{왜 믿을 수 있는지}. 이 보고서의 핵심 질문.
\item \textbf{6장:} 기대수익 점수와 IRR의 역할, 왜 ``아무거나 고를 때''와 비교하는지, Sharpe 정의 세 가지.
\item \textbf{7장:} 결과. 두 축, 등급 안 정보, 등급별 위험, 단순 규칙.
\item \textbf{8장:} 검증과 정정의 기록. \textbf{9장:} 한계와 test 계획.
\end{itemize}
각 핵심 박스 끝의 \textbf{직관:} 줄만 순서대로 읽어도 흐름이 잡힌다.
